\Chapter{Rätzel}
%\addthumb{Rätzel}{\space\Huge\sffamily\bfseries \thechapter}{white}{bbe}
\newpage

%%%%%%%%%%%%%%%%%%%%%%%%%%%%%%%%%%%%%%%%%%%%%%%%%%%%%%%%%%%%%%%%%%%%%%%%%%%%%%%%%%%%%%%%%%%%%%%%%%%%
%Dreiecke
%%%%%%%%%%%%%%%%%%%%%%%%%%%%%%%%%%%%%%%%%%%%%%%%%%%%%%%%%%%%%%%%%%%%%%%%%%%%%%%%%%%%%%%%%%%%%%%%%%%%
\Section{Dreiecke }

    \textbf{Ablauf:}\\  Der Spielleiter zieht wiederholt "Dreieke" zwischen den Mitspielern/Gegenständen/Punkte im Raum.\\
    \textbf{Beispiel:}\\ Ich ziehe ein Dreieck von Markus über Eva zu Stefan. Die Mitspieler müssen herausfinden wemm das Dreieck gehört.\\
    \textbf{Lösung:}\\ Es gehört dem das Dreieck der als erstes, nach dem das Dreieck komplet gezogen ist, spricht.\\



%%%%%%%%%%%%%%%%%%%%%%%%%%%%%%%%%%%%%%%%%%%%%%%%%%%%%%%%%%%%%%%%%%%%%%%%%%%%%%%%%%%%%%%%%%%%%%%%%%%%
%I can see the moon in the spoon
%%%%%%%%%%%%%%%%%%%%%%%%%%%%%%%%%%%%%%%%%%%%%%%%%%%%%%%%%%%%%%%%%%%%%%%%%%%%%%%%%%%%%%%%%%%%%%%%%%%%
\Section{I can see the moon in the spoon }
\textbf{Ablauf}
 Die Teilnehmer sitzen (typischerweise) im Kreis. Einer der Teilnehmer hat einen Löffel. Er betrachtet ihn intensiv, und spricht dann folgenden Satz:

I can see the moon in the spoon. (deutsch: Ich kann den Mond im Löffel sehen).

Anschließend gibt er den Löffel weiter. Diejenigen, die das Spiel bereits kennen, geben dann bekannt, ob er den Mond wirklich gesehen hat, oder ob diese eine Lüge ist. \\
\textbf{Lösung}\\
Wichtig ist, ob dem Übergeber des Löffels gedankt wurde! Nur wer sich ordentlich bedankt, wenn er den Löffel erhält, wird auch den Mond im Löffel sehen können. Was die Gruppe dabei als Bedanken akzeptiert, ist ihr überlassen, sollte aber (möglichst) eindeutig sein.\\


%%%%%%%%%%%%%%%%%%%%%%%%%%%%%%%%%%%%%%%%%%%%%%%%%%%%%%%%%%%%%%%%%%%%%%%%%%%%%%%%%%%%%%%%%%%%%%%%%%%%
%Irrenhaus
%%%%%%%%%%%%%%%%%%%%%%%%%%%%%%%%%%%%%%%%%%%%%%%%%%%%%%%%%%%%%%%%%%%%%%%%%%%%%%%%%%%%%%%%%%%%%%%%%%%%
\Section{Irrenhaus }


\textbf{Ablauf:}\\Die nicht eingeweihten Spieler sind die Irrenärzte, die bereits wissenden Spieler stellen die Irren dar. Die Ärzte dürfen Fragen stellen, jeweils ein Irrer antwortet drauf. Die Antwort klingt dabei meist so, als wäre sie Schwachsinn. Tatsächlich steckt aber System dahinter! Jeder, der beweisen kann, das System zu kennen, darf auch zum Irren werden. Natürlich wird die Lösung nicht erraten, sondern die Spieler müssen durch aktives Tun beweisen, das Wissen zu besitzen.\\
\textbf{Lösung}Die Irren beantworten jeweils die vorangegangene Frage (oder, bei sehr kleveren Spielern, noch weiter zurück liegende Fragen). Die erste Frage wird mit einer möglichst lustigen und unsinnigen Antwort beantwortet. \newline
\textbf{Varianten:}
Wie heißt du? → Raumschiff Enterprise (zufällige Antwort) \newline
Wie viel ist drei mal vier $\rightarrow$ Christian (Antwort auf Frage 1) \newline
Wie heißt die Hauptstadt von Deutschland $\rightarrow$ zwölf (Antwort auf Frage 2) \newline
und so weiter \newline
Alternativ können auch andere Systeme benutzt werden, beispielsweise Antworten mit nur einem Wort oder beginnend mit dem jeweils nächsten Buchstaben des Alphabets, etc.




%%%%%%%%%%%%%%%%%%%%%%%%%%%%%%%%%%%%%%%%%%%%%%%%%%%%%%%%%%%%%%%%%%%%%%%%%%%%%%%%%%%%%%%%%%%%%%%%%%%%
%Flasche auf, Flasche zu
%%%%%%%%%%%%%%%%%%%%%%%%%%%%%%%%%%%%%%%%%%%%%%%%%%%%%%%%%%%%%%%%%%%%%%%%%%%%%%%%%%%%%%%%%%%%%%%%%%%%
\Section{Flasche auf, Flasche zu }

    \textbf{Ablauf:}\\Die Teilnehmer sitzen in eine Kreis und dabei wir eine Flasche rumgegeben. Dabei sagt der die sie Weitgibt ob sie offen oder zu. Die die das Spiel schon kennen sagen dann ob die aussage richtig wahr.  \\
   \textbf{Lösung:}\\Die Flasche ist nicht wirchtig, mann kann mit ihr alles machen was man will. Wichtig ist ob die Person nach dem sie gesagt hat ob sie offen ist oder geschloßen den Mund auf oder geschloßen hatt. Offener Mund $\rightarrow$ offene Flasche.\\


%%%%%%%%%%%%%%%%%%%%%%%%%%%%%%%%%%%%%%%%%%%%%%%%%%%%%%%%%%%%%%%%%%%%%%%%%%%%%%%%%%%%%%%%%%%%%%%%%%%%
%Drache töten
%%%%%%%%%%%%%%%%%%%%%%%%%%%%%%%%%%%%%%%%%%%%%%%%%%%%%%%%%%%%%%%%%%%%%%%%%%%%%%%%%%%%%%%%%%%%%%%%%%%%
\Section{Drache töten }

\textbf{Ablauf:} \\ Verraten sollten die Spieler, welche die Lösung für dieses Ratespiel gefunden haben, diese natürlich nicht, um den Spielspaß möglichst lange aufrecht zu halten. Denn die Teilnehmer sitzen zu Beginn des Spiels im Kreis um den Spielleiter, der als einziger die Lösung kennt. Die Spieler müssen nämlich erraten, wie sich ein Drache am besten töten lässt.
\\
\textbf{Lösung:}\\Was sie zu Beginn des Spiels nicht wissen: Es kommt auf den Anfangsbuchstaben des Gegenstandes, den sie nennen, an. Denkbar sind also auch scheinbar unsinnige Kombinationen wie D-egen, R-egenschirm, -A-ugapfel, C-haos, H-unger, E-mmentaler, D-ynamit. Nach einem Fehlversuch beginnt die Raterunde von Neuem. Das Spiel dauert so lange, bis ausreichend Spieler das System durchschaut haben, mit dem sie den Drachen töten können.\\

 Weil es jüngeren Teilnehmern wohl etwas schwer fallen dürfte, das System zu durchschauen, eignet sich dieses Spiel vor allem für Gruppen gemischten Alters oder mit älteren Teilnehmern. Sie lernen dabei auf eine lustige Art eine Fähigkeit, die ihnen im späteren Leben von großem Nutzen sein kann, nämlich Kombinationsgabe und auch die Fähigkeit, abstrakt zu denken. Schließlich haben Begriffe wie beispielsweise Augapfel oder Hunger rein gar nichts mit den Rittersagen aus dem Mittelalter, in welchen stolze Recken furchterregende Drachen getötet haben, zu tun. Trotzdem tragen diese Begriffe zur Lösung des Rätsels bei.\\

    Die Teilnehmer müssen also erst einmal das System durchschauen und anschließend die richtigen Begriffe finden, wobei Letzteres dann vergleichsweise einsam ist. Sind jüngere Teilnehmer in der Runde, die aufgrund ihrer geistigen Entwicklung noch nicht in der Lage sind, das System, welches hinter diesem Spiel steckt, zu durchschauen, sollte der Spielleiter ihnen dies auf kindgerechte Weise vermitteln. Er kann die Erklärung aber unter Umständen auch so vermitteln, dass die Teilnehmer von selbst auf des Rätsels Lösung kommen.


%%%%%%%%%%%%%%%%%%%%%%%%%%%%%%%%%%%%%%%%%%%%%%%%%%%%%%%%%%%%%%%%%%%%%%%%%%%%%%%%%%%%%%%%%%%%%%%%%%%%
%Kaufhaus
%%%%%%%%%%%%%%%%%%%%%%%%%%%%%%%%%%%%%%%%%%%%%%%%%%%%%%%%%%%%%%%%%%%%%%%%%%%%%%%%%%%%%%%%%%%%%%%%%%%%
\Section{Kaufhaus }

\textbf{Ablauf:}\\
Der Spielleiter (der die Lösung des Rätsels kennt) beginnt mit folgender kurzen Geschichte:\\
\textit{Ich gehe ins Kaufhaus.}\\
\textit{Dort gehe ich in den (X) Stock (links oder rechts)}\\
\textit{Ich kaufe mir dort ein (Gegenstand)}\\ 
Danach versuchen die andern Spieler selber etwas in Kaufhaus zu kaufen.\\
\newline
\textbf{Lösung:}\\
Beim Betreten des Kaufhauses geht der fragende Spieler nach links oder rechts; dies beschreibt die Spieler, die links oder rechts neben ihm stehen. Der "2. Stock links" beschreibt also den zweiten Spieler, der links neben dem Fragesteller steht (aus seiner Sicht). Geht der Spieler ins Erdgeschoss ("0. Stock"), ist der Fragesteller selber gemeint. Der Spieler der dort sitzt muss den Gegenstand an sich haben.

%%%%%%%%%%%%%%%%%%%%%%%%%%%%%%%%%%%%%%%%%%%%%%%%%%%%%%%%%%%%%%%%%%%%%%%%%%%%%%%%%%%%%%%%%%%%%%%%%%%%
%Türsteher
%%%%%%%%%%%%%%%%%%%%%%%%%%%%%%%%%%%%%%%%%%%%%%%%%%%%%%%%%%%%%%%%%%%%%%%%%%%%%%%%%%%%%%%%%%%%%%%%%%%%
\Section{Türsteher }
\textbf{Beispiel:}\\

\textit{Ich gehe in den Club und zeige bringe den Türsteher eine Bild von meiner Oma mit und ich komme rein.}\\
\textbf{Lösung}\\
Wie genau der Satz gesagt wird ist nicht wichtig, nur dass eine Zeitung/Zeitschrift drin vorkommt. Was gesagt werden muss ist: \textit{Ich gehe in den Club} und \textit{und ich komme rein.}

%%%%%%%%%%%%%%%%%%%%%%%%%%%%%%%%%%%%%%%%%%%%%%%%%%%%%%%%%%%%%%%%%%%%%%%%%%%%%%%%%%%%%%%%%%%%%%%%%%%%
%Insel keinIkeinU
%%%%%%%%%%%%%%%%%%%%%%%%%%%%%%%%%%%%%%%%%%%%%%%%%%%%%%%%%%%%%%%%%%%%%%%%%%%%%%%%%%%%%%%%%%%%%%%%%%%%
\Section{Insel Keinikeinu }
Bei diesem Spiel verreisen die Mitspieler:innen in das Keinikeinuland.
Sie packen ihren Koffer und dürfen einige Sachen mitnehmen, allerdings dürfen sie nicht alles mitnehmen, bestimmte Dinge scheiden aus.
Wer das Prinzip erkannt hat, darf mit ins Keinikeinuland.
Die Mitspieler:innen sagen der Reihe nach, was sie mitnehmen möchten, die Spielleitung und die Spieler, die es wissen, antworten, ob er oder sie mit ihnen fahren darf oder nicht.

\textbf{Lösung:}
Die mitgenommenen Sachen dürfen, wie der Name des Spiels schon sagt, kein ‘i’ und kein ‘u’ enthalten.




%%%%%%%%%%%%%%%%%%%%%%%%%%%%%%%%%%%%%%%%%%%%%%%%%%%%%%%%%%%%%%%%%%%%%%%%%%%%%%%%%%%%%%%%%%%%%%%%%%%%
%Maus-Maus-Elefant
%%%%%%%%%%%%%%%%%%%%%%%%%%%%%%%%%%%%%%%%%%%%%%%%%%%%%%%%%%%%%%%%%%%%%%%%%%%%%%%%%%%%%%%%%%%%%%%%%%%%
\Section{Maus-Maus-Elefant }
Affe-Affe-Elefant ist ein Rätsel. Ein Teilnehmer zeigt eine Bewegungsabfolge mit Text vor, die anderen müssen sie exakt nachmachen. Klingt einfacher, als es ist!\\
\textbf{Ablauf:}\\
Ein Spieler zeigt eine Bewegungsabfolge mit den Händen vor und spricht einen ebenso einfachen wie sinnlosen Text dazu.
\begin{table}[h]
    \begin{tabular}{|l|l|}
    \hline
         \textbf{Text} & \textbf{Bewegung}  \\ \hline 
         Also: & \\ \hline
         Maus & tippt auf kleinen Finger \\ \hline
         Maus   &        tippt auf den Ringfinger \\ \hline
         Maus & tippt auf den Mittelfinger \\ \hline
         Maus & tippt auf den Zeigefinger \\ \hline
         Elefant & streicht den Zeigefinger entlang nach unten und am Daumen wieder hoch \\ \hline
         Maus & tippt auf den Daumen \\ \hline
         Elefant & streicht zurück über Daumen und Zeigefinger \\ \hline
         Maus & tippt auf den Zeigefinger\\ \hline
         Maus & tippt auf den Mittelfinger \\ \hline
         Maus & tippt auf den Ringfinger \\ \hline
         Maus & tippt auf kleinen Finger \\ \hline
         Und jetzt mach's nach! & \\ \hline
         
    \end{tabular}
\end{table}

